\chapter{Construire une application, de zéro à l'exécutable}

Nous avons vu la pile (chapitre~1), le rendu (chapitre~2) et les widgets
(chapitre~3). Il reste à tout assembler. Ce chapitre construit une application
complète, dans l'ordre, en expliquant à chaque étape \emph{pourquoi} elle est là
et \emph{ce qui casse} si on l'oublie.

À la fin, vous aurez un programme qui compile, se lance, affiche une interface et
répond à la souris et au clavier.

\section{Vue d'ensemble du squelette}

Avant le détail, voici la carte. Une application NKGui sur NkCanvas se construit
en sept étapes, puis boucle.

\begin{nkcode}{Les sept étapes de l'initialisation}
1. FENETRE          NkWindow + NkWindowConfig
2. CONTEXTE GPU     NkContextDesc  (choix de l'API graphique)
3. CIBLE DE RENDU   NkRenderWindow (possede le cycle de frame)
4. CONTEXTE UI      NkGuiContext::Init + SetCurrentContext
5. BACKEND (PONT)   NkGuiCanvasBackend::Init(target->GetRenderer())
6. POLICE           NkGuiFont::LoadEmbedded  PUIS  backend.UploadFontGray8
7. ENTREE           callbacks NKEvent -> ctx.input

BOUCLE :
   evenements -> [resize] -> ctx.BeginFrame(dt) -> widgets -> ctx.EndFrame()
   -> target->Clear() -> backend.Submit(ctx.dl) -> backend.Submit(ctx.dlOverlay)
   -> target->Display()
\end{nkcode}

L'ordre n'est pas indifférent. Le pont a besoin du renderer, donc de la cible.
L'upload de l'atlas a besoin du pont \emph{et} de la police chargée. Les
\emph{callbacks} ont besoin du contexte, puisqu'ils écrivent dedans.

\section{Étape 1 — la fenêtre}

\begin{nkcode}{Applications/NKGuiDemo/src/NKGuiDemo/main.cpp:243-251}
    NkWindow window;
    NkWindowConfig cfg;
    cfg.title = "NKGui - Demo Phase 2 (Coeur : drawlist + interaction)";
    cfg.width = 1100;
    cfg.height = 800;
    cfg.centered = true;
    cfg.resizable = true;
    if (!window.Create(cfg))
        return -1;
\end{nkcode}

Rien de subtil, mais notez le test de retour : \texttt{Create} renvoie
\texttt{false} si la fenêtre n'a pas pu être créée, et il n'y a pas d'exception à
attraper. Chaque étape de cette initialisation suit le même modèle — retour
booléen ou méthode \texttt{IsValid()} — et chacune doit être testée.

Les en-têtes nécessaires :

\begin{nkcode}{Applications/NKGuiDemo/src/NKGuiDemo/main.cpp:7-11 (extrait)}
#include "NKWindow/NKWindow.h"
#include "NKWindow/NKMain.h"
#include "NKEvent/NkWindowEvent.h"
#include "NKEvent/NkMouseEvent.h"
#include "NKEvent/NkKeyboardEvent.h"
\end{nkcode}

\section{Étape 2 et 3 — contexte GPU et cible de rendu}

\begin{nkcode}{Applications/NKGuiDemo/src/NKGuiDemo/main.cpp:253-265}
    NkContextDesc desc;
    desc.api = NkGraphicsApi::NK_GFX_API_AUTO;
    if (desc.api == NkGraphicsApi::NK_GFX_API_AUTO) {
#if defined(NKENTSEU_PLATFORM_WINDOWS)
        desc.api = NkGraphicsApi::NK_GFX_API_DX11;
#else
        desc.api = NkGraphicsApi::NK_GFX_API_OPENGL;
#endif
    }
    auto target = memory::NkMakeUnique<NkRenderWindow>(window, desc);
    if (!target || !target->IsValid())
        return -1;
\end{nkcode}

Trois observations :

\begin{itemize}
  \item \texttt{NK\_GFX\_API\_AUTO} est résolu ici, à la main, par une directive
        de préprocesseur. On aurait pu utiliser
        \texttt{NkContextFactory::CreateWithFallback} et laisser le moteur
        essayer une liste ordonnée ; la démo préfère être explicite.
  \item La cible est créée par \texttt{memory::NkMakeUnique} et vivra jusqu'à la
        fin de \texttt{nkmain}. Elle possède le contexte graphique \emph{et} le
        renderer 2D.
  \item Le double test \texttt{!target || !target->IsValid()} couvre les deux
        modes d'échec : l'allocation, et l'initialisation GPU.
\end{itemize}

Les en-têtes correspondants :

\begin{nkcode}{Applications/NKGuiDemo/src/NKGuiDemo/main.cpp:12-16 (extrait)}
#include "NKCanvas/Core/NkContextDesc.h"
#include "NKCanvas/Core/NkGraphicsApi.h"
#include "NKCanvas/Renderer/Targets/NkRenderWindow.h"
#include "NKTime/NkClock.h"
#include "NKMemory/NkUniquePtr.h"
\end{nkcode}

\begin{nkretenir}
Si votre fenêtre reste noire alors que le programme tourne, changez d'API
graphique avant de chercher dans votre code. Rappel du chapitre~1 : seul OpenGL
était validé à l'exécution au 30 mai 2026 selon
\nkfile{Kernel/Runtime/NKCanvas/ROADMAP.md}, et DX11 comme le rasteriseur
logiciel affichaient un écran vide. Le journal
(\nkfile{logs/app.log}) vous dit quel backend a réellement été retenu.
\end{nkretenir}

\section{Étape 4 — le contexte NKGui}

\begin{nkcode}{Applications/NKGuiDemo/src/NKGuiDemo/main.cpp:266-272}
    auto ctxPtr = memory::NkMakeUnique<NkGuiContext>();
    if (!ctxPtr)
        return -1;
    NkGuiContext &ctx = *ctxPtr;
    ctx.Init(static_cast<int32>(cfg.width), static_cast<int32>(cfg.height));
    ctx.styleFn = &DemoStyle; // hook de style (override de dessin par widget)
    SetCurrentContext(&ctx);
\end{nkcode}

Le contexte est volumineux — plus de cent champs et trente-deux listes de dessin
de fenêtres — ce qui explique qu'on l'alloue plutôt que de le poser sur la pile.
On prend ensuite une référence, ce qui rend tout le code des widgets lisible.

\texttt{Init(width, height)} pose la taille de vue initiale et prépare les
structures. \texttt{SetCurrentContext} installe le contexte courant du fil
d'exécution ; la ligne \texttt{ctx.styleFn} est facultative.

C'est ici, juste après \texttt{Init}, que vous personnaliserez le thème :

\begin{nkcode}{Personnalisation du thème — écrit pour ce cours, d'après NkEditorShell.cpp:155-179}
    ctx.Init(1100, 800);
    ctx.theme.bgPrimary = {13, 17, 23, 255};
    ctx.theme.panel     = {22, 27, 34, 255};
    ctx.theme.accent    = {88, 166, 255, 255};
    ctx.theme.rounding  = 6.f;
    SetCurrentContext(&ctx);
\end{nkcode}

\section{Étape 5 — le pont vers NkCanvas}

\begin{nkcode}{Applications/NKGuiDemo/src/NKGuiDemo/main.cpp:274-276}
    renderer::NkGuiCanvasBackend backend;
    if (!backend.Init(target->GetRenderer()))
        return -1;
\end{nkcode}

Deux lignes, et c'est tout le mariage entre les deux modules. L'en-tête à inclure
est celui qui vit dans NkCanvas :

\begin{nkcode}{Applications/NKGuiDemo/src/NKGuiDemo/main.cpp:25}
#include "NKCanvas/UI/NkGuiCanvasBackend.h" // backend NKGui->NKCanvas (lib reutilisable)
\end{nkcode}

Rappelez-vous du chapitre~1 : ce fichier est \emph{header-only}. Il n'est compilé
nulle part ailleurs que dans votre propre unité de traduction. C'est pour cela
que votre \texttt{.jenga} doit dépendre à la fois de \texttt{NKGui} et de
\texttt{NKCanvas}.

\texttt{backend} est déclaré sur la pile de \texttt{nkmain}. Son destructeur
libère les textures qu'il a créées :

\begin{nkcode}{Kernel/Runtime/NKCanvas/src/NKCanvas/UI/NkGuiCanvasBackend.h:24-32}
~NkGuiCanvasBackend() {
    auto &alloc = nkentseu::memory::NkGetDefaultAllocator();
    for (uint32 i = 0; i < mFonts.Size();  ++i) if (mFonts[i].tex)  alloc.Delete(mFonts[i].tex);
    for (uint32 i = 0; i < mImages.Size(); ++i) if (mImages[i].tex) alloc.Delete(mImages[i].tex);
}
\end{nkcode}

\begin{nkpiege}[L'ordre de destruction]
Le pont détruit des textures GPU, donc il doit mourir \emph{avant} le renderer
qui les a créées. Comme \texttt{backend} est déclaré \emph{après} \texttt{target}
dans \texttt{nkmain}, il est détruit \emph{avant} — les objets locaux se
détruisent dans l'ordre inverse de leur déclaration. Cet ordre est correct par
construction ; inverser les deux déclarations produirait une libération de
textures sur un renderer déjà mort.
\end{nkpiege}

\section{Étape 6 — la police, et l'upload de son atlas}

C'est l'étape où l'on se trompe le plus souvent, parce qu'elle est en deux temps
et que rien ne signale l'oubli du second.

\begin{nkcode}{Applications/NKGuiDemo/src/NKGuiDemo/main.cpp:281-290}
    auto fontPtr = memory::NkMakeUnique<NkGuiFont>();
    if (!fontPtr->LoadEmbedded(NkEmbeddedFontId::DroidSans, 18.f)) {
        fontPtr->LoadEmbedded(NkEmbeddedFontId::ProggyClean, 16.f);
    }
    ctx.font = fontPtr.Get();
    if (fontPtr->Valid()) {
        backend.UploadFontGray8(fontPtr->TexId(), fontPtr->pixels, fontPtr->atlasW, fontPtr->atlasH);
    }
\end{nkcode}

\subsection{Ce que fait \texttt{LoadEmbedded}}

\begin{nkcode}{Kernel/Runtime/NKGui/src/NKGui/Core/NkGuiFont.h:19-40 (abrégé)}
struct NKENTSEU_NKGUI_CLASS_EXPORT NkGuiFont {
        NkFontAtlas atlas;          ///< possède la texture + les glyphes
        NkFont *face = nullptr;     ///< face produite (détenue par l'atlas)
        uint32 texId = 0x4E4B4654u; ///< 'NKFT' — id stable pour le backend
        uint8 *pixels = nullptr;    ///< atlas alpha8 (détenu par l'atlas)
        int32 atlasW = 0;
        int32 atlasH = 0;
        bool dirty = false;         ///< à (ré)uploader côté backend

        NkGuiFont() = default;
        NkGuiFont(const NkGuiFont &) = delete;            // NON COPIABLE
        NkGuiFont &operator=(const NkGuiFont &) = delete;

        bool LoadEmbedded(NkEmbeddedFontId id, float32 sizePx, bool extFallback = true) noexcept;
        bool LoadFromFile(const char *path, float32 sizePx, bool extFallback = true) noexcept;
\end{nkcode}

L'appel rasterise les glyphes dans un atlas et laisse \texttt{pixels} pointer sur
un bitmap en niveaux de gris (un octet par pixel, l'alpha). \textbf{Cet atlas est
en mémoire centrale ; le GPU ne le connaît pas encore.}

Les polices embarquées disponibles
(\nkfile{Kernel/Runtime/NKFont/src/NKFont/Embedded/NkFontEmbedded.h:49-67}) :
\texttt{ProggyClean}, \texttt{ProggyTiny} (bitmap monospace),
\texttt{DroidSans}, \texttt{Karla}, \texttt{Roboto} (sans-serif),
\texttt{Cousine}, \texttt{SourceCodePro} (monospace vectorielles),
\texttt{DroidSerif}, \texttt{DejaVuSansMono}.

\subsection{Ce que fait \texttt{UploadFontGray8}}

Le pont convertit l'atlas en RGBA — blanc partout, l'alpha portant la forme du
glyphe — et crée la texture GPU :

\begin{nkcode}{Kernel/Runtime/NKCanvas/src/NKCanvas/UI/NkGuiCanvasBackend.h:47-55}
const usize n = (usize)w * (usize)h;
mExpand.Resize(n * 4u);
uint8 *d = mExpand.Data();
for (usize i = 0; i < n; ++i) {
    d[i*4+0] = 255u; d[i*4+1] = 255u; d[i*4+2] = 255u; d[i*4+3] = gray[i];
}
\end{nkcode}

C'est là que le \texttt{texId} prend son sens : le pont range la texture dans une
table indexée par cet identifiant, et \texttt{Submit} la retrouvera quand une
commande de dessin portera le même. La constante par défaut est
\texttt{0x4E4B4654} — les quatre lettres « NKFT ». Elle n'est \textbf{jamais
nulle}, car zéro est réservé à la texture blanche interne côté RHI.

\begin{nkpiege}[Oublier l'upload : du texte parfaitement invisible]
Si vous chargez la police et posez \texttt{ctx.font} mais oubliez
\texttt{UploadFontGray8}, il ne se passe \emph{rien} de visible et
\emph{aucun} message n'apparaît. Les widgets se dessinent (leurs rectangles sont
là), le layout est correct — mais chaque commande de texte référence un
\texttt{texId} que le pont ne connaît pas, et elle est silencieusement ignorée.

Symptôme : des boutons vides, des cases à cocher sans étiquette, un panneau au
titre absent. Diagnostic : vérifiez que \texttt{fontPtr->Valid()} est vrai, et
que l'upload est bien appelé après le chargement.
\end{nkpiege}

\begin{nkpiege}[La police ne doit jamais mourir avant la boucle]
\begin{nkcode}{Kernel/Runtime/NKGui/src/NKGui/Core/NkGuiContext.h:141-142 (sens)}
NkGuiFont *font = nullptr;      ///< pointeur BRUT, NON POSSEDE
NkGuiFont *codeFont = nullptr;  ///< idem
\end{nkcode}

\texttt{ctx.font = fontPtr.Get()} donne au contexte un pointeur \emph{brut}. Le
contexte ne possède rien. Si l'objet \texttt{NkGuiFont} est détruit — parce qu'il
était local à une fonction d'initialisation, par exemple — le contexte pointe sur
de la mémoire libérée et la première mesure de texte plante.

La règle : \textbf{l'objet police doit vivre au moins aussi longtemps que la
boucle}. La démo le tient dans un \texttt{NkUniquePtr} déclaré dans
\texttt{nkmain}, à côté de tout le reste.
\end{nkpiege}

\begin{nkpiege}[Recharger un atlas à une autre taille]
\begin{nkcode}{Kernel/Runtime/NKCanvas/src/NKCanvas/UI/NkGuiCanvasBackend.h:65-66}
// (Re)créer la texture si elle n'existe pas OU si la taille change (rechargement
// de police à une autre taille = zoom / DPI). Sinon Update() déborderait l'ancienne taille.
\end{nkcode}
Le pont gère ce cas pour vous. Mais \textbf{vous} devez veiller à ne jamais
recharger une police au milieu d'une image : les commandes déjà émises
référenceraient un atlas remplacé, et leurs coordonnées de texture ne
correspondraient plus. Faites-le avant \texttt{BeginFrame}, comme la démo le fait
pour la touche F9.
\end{nkpiege}

\subsection{Les polices de repli}

Pour les alphabets que la police principale ne couvre pas :

\begin{nkcode}{Kernel/Runtime/NKGui/src/NKGui/Core/NkGuiFont.h:71-78}
// Polices de REPLI EXTERNES (fichiers .ttf charges au runtime) : tout glyphe
// absent des polices principales (Inter/DejaVu) y est cherche. Roles :
//   broad : large couverture (latin etendu, grec, cyrillique, hebreu, arabe, symboles)
//   cjk   : ideogrammes (chinois/japonais/coreen) - volumineux, opt-in
//   emoji : emoji monochrome
// A poser par l'APPLICATION (ex. NKCode, depuis son dossier data/fonts) AVANT
// de charger les polices. Un chemin nullptr/vide = role desactive.
void NkSetFallbackFontPaths(const char *broad, const char *cjk, const char *emoji) noexcept;
\end{nkcode}

\textbf{Avant} de charger la police, donc. C'est un invariant d'ordre : l'appel
place des chemins dans des variables globales que la construction de l'atlas
consultera.

Pour une police monospace destinée au code, on désactive au contraire les replis
(\texttt{extFallback = false}) : « atlas minuscule => reconstruction RAPIDE au
zoom » (\nkfile{NkGuiFont.h:34-35}).

\section{Étape 7 — brancher l'entrée}

Les \emph{callbacks} écrivent directement dans \texttt{ctx.input}. Voici
l'ensemble minimal, puis les compléments.

\subsection{Souris et fermeture}

\begin{nkcode}{Applications/NKGuiDemo/src/NKGuiDemo/main.cpp:326-343 (extrait)}
    auto &events = NkEvents();
    bool running = true;
    events.AddEventCallback<NkWindowCloseEvent>([&](NkWindowCloseEvent *) { running = false; });

    events.AddEventCallback<NkMouseMoveEvent>([&](NkMouseMoveEvent *e) {
        ctx.input.mousePos = {static_cast<float32>(e->GetX()), static_cast<float32>(e->GetY())};
    });
    events.AddEventCallback<NkMouseButtonPressEvent>([&](NkMouseButtonPressEvent *e) {
        if (e->GetButton() == NkMouseButton::NK_MB_LEFT)  ctx.input.mouseDown[0] = true;
        if (e->GetButton() == NkMouseButton::NK_MB_RIGHT) ctx.input.mouseDown[1] = true;
        ctx.input.ctrlDown  = e->GetModifiers().ctrl;
        ctx.input.shiftDown = e->GetModifiers().shift;
        ctx.input.altDown   = e->GetModifiers().alt;
    });
    events.AddEventCallback<NkMouseButtonReleaseEvent>([&](NkMouseButtonReleaseEvent *e) {
        if (e->GetButton() == NkMouseButton::NK_MB_LEFT)  ctx.input.mouseDown[0] = false;
        if (e->GetButton() == NkMouseButton::NK_MB_RIGHT) ctx.input.mouseDown[1] = false;
    });
\end{nkcode}

Indices des boutons : \texttt{0} = gauche, \texttt{1} = droit, \texttt{2} =
milieu.

\subsection{La molette}

\begin{nkcode}{Applications/NKGuiDemo/src/NKGuiDemo/main.cpp:344-352 (extrait)}
    events.AddEventCallback<NkMouseWheelVerticalEvent>([&](NkMouseWheelVerticalEvent *e) {
        ctx.input.wheel += static_cast<float32>(e->GetDeltaY()); // >0 = haut ; consommé en EndFrame
        const auto m = e->GetModifiers();                        // modificateurs au moment de la molette
        ctx.input.ctrlDown = m.ctrl; ctx.input.shiftDown = m.shift; ctx.input.altDown = m.alt;
    });
    events.AddEventCallback<NkMouseWheelHorizontalEvent>([&](NkMouseWheelHorizontalEvent *e) {
        ctx.input.wheelH += static_cast<float32>(e->GetDeltaX());
    });
\end{nkcode}

Notez le \texttt{+=} et non le \texttt{=} : plusieurs crans peuvent arriver dans
la même image, ils s'accumulent. \texttt{EndFrame} remet le compteur à zéro.
Notez aussi la relecture des modificateurs : sans elle, le Ctrl+molette (zoom)
ne fonctionnerait pas de façon fiable.

\subsection{Le double-clic}

\begin{nkcode}{Applications/NKGuiDemo/src/NKGuiDemo/main.cpp:353-359 (extrait)}
    // Double-clic détecté par l'OS : NKEvent l'émet comme événement DÉDIÉ (le 2e
    // clic n'arrive PAS comme un mouseDown normal) → on l'injecte explicitement.
    events.AddEventCallback<NkMouseDoubleClickEvent>([&](NkMouseDoubleClickEvent *e) {
        ctx.input.mousePos = {static_cast<float32>(e->GetX()), static_cast<float32>(e->GetY())};
        if (e->GetButton() == NkMouseButton::NK_MB_LEFT) ctx.input.SetDoubleClick(0);
    });
\end{nkcode}

Sans ce \emph{callback}, le second clic d'un double-clic n'existe tout simplement
pas pour NKGui : le système l'a converti en événement dédié. Les renommages en
place et la saisie au clavier dans les \texttt{DragFloat}, qui reposent sur le
double-clic, ne fonctionneraient pas.

\subsection{Le texte et les touches}

\begin{nkcode}{Applications/NKGuiDemo/src/NKGuiDemo/main.cpp:360-380 (extrait)}
    // Saisie texte + touches d'édition (pour InputText). On pose l'état ENFONCÉ
    // (press/release) pour permettre la répétition au maintien.
    events.AddEventCallback<NkTextInputEvent>([&](NkTextInputEvent *e) { ctx.input.PushChar(e->GetCodepoint()); });

    auto setKey = [&](NkKey k, bool down) {
        switch (k) {
            case NkKey::NK_LEFT:   ctx.input.SetKey(NkGuiKey::Left, down);      break;
            case NkKey::NK_RIGHT:  ctx.input.SetKey(NkGuiKey::Right, down);     break;
            case NkKey::NK_UP:     ctx.input.SetKey(NkGuiKey::Up, down);        break;
            case NkKey::NK_DOWN:   ctx.input.SetKey(NkGuiKey::Down, down);      break;
            case NkKey::NK_HOME:   ctx.input.SetKey(NkGuiKey::Home, down);      break;
            case NkKey::NK_END:    ctx.input.SetKey(NkGuiKey::End, down);       break;
            case NkKey::NK_BACK:   ctx.input.SetKey(NkGuiKey::Backspace, down); break;
            case NkKey::NK_DELETE: ctx.input.SetKey(NkGuiKey::Delete, down);    break;
            case NkKey::NK_ENTER:  ctx.input.SetKey(NkGuiKey::Enter, down);     break;
            case NkKey::NK_ESCAPE: ctx.input.SetKey(NkGuiKey::Escape, down);    break;
            default: break;
        }
    };
    events.AddEventCallback<NkKeyPressEvent>([&](NkKeyPressEvent *e)    { setKey(e->GetKey(), true);  … });
    events.AddEventCallback<NkKeyReleaseEvent>([&](NkKeyReleaseEvent *e){ setKey(e->GetKey(), false); … });
\end{nkcode}

Deux points essentiels, déjà annoncés au chapitre~1 :

\begin{itemize}
  \item on pose l'état \textbf{enfoncé/relâché}, pas un « appui » ponctuel. C'est
        ce qui permet à NKGui de calculer les durées et donc la répétition au
        maintien ;
  \item le \textbf{texte} passe par \texttt{PushChar} et les \textbf{touches} par
        \texttt{SetKey}. Ce sont deux canaux séparés.
\end{itemize}

Il reste à brancher les raccourcis d'édition, sur le modèle du shell de
l'éditeur :

\begin{nkcode}{Engine/NKEditorKit/src/NKEditorKit/NkEditorShell.cpp:318-326 (extrait)}
if (e->GetModifiers().ctrl) { // raccourcis copier/couper/coller/tout-selectionner
    if      (k == NkKey::NK_C) mUI.input.wantCopy = true;
    else if (k == NkKey::NK_X) mUI.input.wantCut = true;
    else if (k == NkKey::NK_V) mUI.input.wantPaste = true;
    else if (k == NkKey::NK_A) mUI.input.wantSelectAll = true;
\end{nkcode}

\subsection{Le presse-papiers}

NKGui ne connaît pas la fenêtre, donc il ne peut pas accéder au presse-papiers du
système. On lui donne deux pointeurs de fonction
(\texttt{ctx.clipboardGetFn}, \texttt{ctx.clipboardSetFn}, plus un pointeur
utilisateur \texttt{ctx.clipboardUser}) ; c'est ce que fait le shell en
\nkfile{NkEditorShell.cpp:185-190}. Sans cela, copier-coller dans un champ de
saisie ne fait rien.

\section{La boucle principale}

\subsection{Le temps}

\begin{nkcode}{Applications/NKGuiDemo/src/NKGuiDemo/main.cpp:431-436}
    while (running && window.IsOpen()) {
        float32 dt = clock.Tick().delta;
        if (dt <= 0.f)
            dt = 1.f / 60.f;
        if (dt > 0.1f)
            dt = 0.1f;
\end{nkcode}

Les deux bornes ont chacune leur raison : \texttt{dt <= 0} arrive à la toute
première image (l'horloge n'a pas encore de référence) ; \texttt{dt > 0.1} arrive
après un blocage — déplacement de fenêtre, mise en veille, point d'arrêt du
débogueur. Sans la borne haute, tout ce qui dépend du temps ferait un bond : le
curseur de saisie clignoterait plusieurs fois d'un coup, les répétitions au
maintien se déclencheraient en rafale.

\subsection{Les événements — avant tout le reste}

\begin{nkcode}{Applications/NKGuiDemo/src/NKGuiDemo/main.cpp:438-442}
        while (NkEvent *ev = NkEvents().PollEvent()) {
            (void)ev;
        }
        if (!running)
            break;
\end{nkcode}

Le \texttt{if (!running) break;} évite de dessiner une image entière sur une
fenêtre qu'on vient de fermer.

\subsection{Le redimensionnement}

\begin{nkcode}{Applications/NKGuiDemo/src/NKGuiDemo/main.cpp:444-458}
        // Synchroniser la SWAPCHAIN à la taille de la FENÊTRE. target->GetSize()
        // renvoie la taille de la swapchain (qui ne change QU'À OnResize) — s'en
        // servir pour détecter le resize était faux : la swapchain restait figée à
        // sa taille de création (rendu basse-résolution étiré). On pilote OnResize
        // depuis la taille fenêtre (inclut la 1re frame → sync initiale).
        const math::NkVec2u wsz = target->GetWindow().GetSize();
        if (wsz.x > 0 && wsz.y > 0 && (wsz.x != lastW || wsz.y != lastH)) {
            target->OnResize(wsz.x, wsz.y);
            lastW = wsz.x;
            lastH = wsz.y;
        }
        const math::NkVec2u sz = target->GetSize(); // = wsz après OnResize
        if (sz.x > 0 && sz.y > 0) {
            ctx.viewW = static_cast<int32>(sz.x);
            ctx.viewH = static_cast<int32>(sz.y);
        }
\end{nkcode}

\texttt{lastW} et \texttt{lastH} sont initialisés à zéro, ce qui garantit que la
première image déclenche une synchronisation. Et les tests \texttt{> 0} évitent
de recréer une chaîne d'échange de taille nulle quand la fenêtre est réduite.

\begin{nkpiege}[Fenêtre minimisée]
Une fenêtre en cours de réduction glisse un rectangle intermédiaire — environ
$160 \times 28$ sous Windows — entre le test « minimisée ? » et la lecture de
taille. Le shell de l'éditeur documente une garde qui refuse toute taille sous
32 pixels :

\begin{nkcode}{Integrations/NKGui/NkEditorRHIRenderer.h:123-130 (extrait)}
// GARDE ANTI-MINIMISATION : une fenetre en cours de reduction
// glisse son rect placeholder (~160x28 sous Windows) entre le
// test « minimisee ? » de la boucle principale et la lecture
// de taille -- la course est reelle (defaut 4.3, reproduite).
// Sous 32 px, les cibles divisees du rendu (bloom /32) tombent
// a zero et CreateTexture2D echoue en E_INVALIDARG -> mort a
// la restauration. On refuse net : la vraie taille arrivera
// avec le retour de la fenetre.
\end{nkcode}
\end{nkpiege}

\subsection{L'image d'interface}

\begin{nkcode}{Applications/NKGuiDemo/src/NKGuiDemo/main.cpp:468-481 (extrait)}
        ctx.BeginFrame(dt);
        NkGuiDrawList &dl = ctx.dl;
        const float32 W = static_cast<float32>(ctx.viewW);
        const float32 H = static_cast<float32>(ctx.viewH);

        // Fond + barre d'en-tête + TITRE (texte NKFont).
        dl.AddRectFilled({0.f, 0.f, W, H}, ctx.theme.bgPrimary);
        dl.AddRectFilled({0.f, 0.f, W, 40.f}, ctx.theme.header);
        dl.AddRectFilled({0.f, 38.f, W, 2.f}, ctx.theme.accent);
        TextAt(ctx, {16.f, 11.f}, "NKGui - Phase 3 : layout + widgets (NKFont)");
\end{nkcode}

Le fond est dessiné en premier — tout le reste passera par-dessus. Ici la démo
écrit dans \texttt{ctx.dl} directement, ce qui est légitime parce qu'aucune
fenêtre n'est ouverte à ce moment : nous sommes bien sur la couche principale.
Dès qu'on est \emph{à l'intérieur} d'un \texttt{Begin}, il faut passer par
\texttt{ctx.DL()} (chapitre~3).

Puis les widgets, avec une région de layout ouverte :

\begin{nkcode}{Squelette de widgets — écrit pour ce cours}
        ctx.BeginLayout({20.f, 60.f, 400.f, H - 80.f});
        Text(ctx, "Bonjour NKGui");
        if (Button(ctx, "Cliquez-moi")) {
            logger.Info("clic !");
        }
        static bool visible = true;
        Checkbox(ctx, "Afficher le panneau", visible);
        static float32 zoom = 1.f;
        SliderFloat(ctx, "Zoom", zoom, 0.5f, 4.f);
\end{nkcode}

\subsection{Le curseur}

\begin{nkcode}{Applications/NKGuiDemo/src/NKGuiDemo/main.cpp:1020-1041}
        // Curseur souhaité par NKGui → curseur OS (OPTIONNEL : l'app choisit d'appliquer).
        // Ex. DragFloat pose ↔ / ↕ pendant le survol/glisser. SetCursor est persistant.
        {
            NkWindow::NkCursorType c = NkWindow::NkCursorType::Arrow;
            switch (ctx.wantCursor) {
                case NkGuiCursor::Text:
                    c = NkWindow::NkCursorType::TextInput;
                    break;
                case NkGuiCursor::Hand:
                    c = NkWindow::NkCursorType::Hand;
                    break;
                case NkGuiCursor::ResizeEW:
                    c = NkWindow::NkCursorType::ResizeWE;
                    break;
                case NkGuiCursor::ResizeNS:
                    c = NkWindow::NkCursorType::ResizeNS;
                    break;
                default:
                    break;
            }
            window.SetCursor(c);
        }
\end{nkcode}

Rappel du chapitre~3 : NKGui \emph{demande} un curseur en posant
\texttt{ctx.wantCursor}, il ne le pose pas — il ne connaît pas la fenêtre. Ce
bloc est le traducteur, et il est facultatif : sans lui, tout fonctionne, mais
le pointeur ne change jamais de forme au-dessus d'un champ de saisie ou d'un
séparateur. Notez qu'il est placé \emph{avant} \texttt{EndFrame}, parce que
\texttt{wantCursor} est réinitialisé par \texttt{BeginFrame} et rempli par les
widgets.

\subsection{La présentation}

\begin{nkcode}{Applications/NKGuiDemo/src/NKGuiDemo/main.cpp:1043-1048}
        ctx.EndFrame();

        target->Clear();
        backend.Submit(dl, sz.x, sz.y);               // couche principale
        backend.Submit(ctx.dlOverlay, sz.x, sz.y); // couche popups/overlay (au-dessus)
        target->Display(); // Capture d'ecran (F12) : APRES Display() — la frame presentee. Self-capture
\end{nkcode}

\begin{nkpiege}[DEUX \texttt{Submit}, jamais un seul]
C'est le piège le plus coûteux de tout ce cours, parce qu'il est absolument
silencieux.

NKGui produit \textbf{deux} listes de dessin :
\begin{itemize}
  \item \texttt{ctx.dl} — le contenu ordinaire, y compris les fenêtres fusionnées
        par ordre de profondeur ;
  \item \texttt{ctx.dlOverlay} — les popups, les menus déroulants, les
        sous-menus, les combos, les infobulles, et tout ce que vous dessinez
        entre \texttt{PushOverlay} et \texttt{PopOverlay}.
\end{itemize}

Si vous n'appelez \texttt{Submit} que sur la première, \textbf{tout fonctionne
sauf les surfaces flottantes}. Le menu s'ouvre — l'état est correct, les clics
sont capturés, la logique tourne — mais rien ne s'affiche. Vous cliquez dans le
vide sur des éléments invisibles.

Et il n'y a \emph{aucun} message d'erreur : le framework a fait son travail, le
pont aussi ; c'est simplement qu'on ne lui a jamais donné la seconde liste.

\textbf{L'ordre compte aussi} : \texttt{ctx.dl} d'abord, \texttt{ctx.dlOverlay}
ensuite. Inversé, les popups passeraient \emph{sous} le contenu.
\end{nkpiege}

Le shell de l'éditeur fait exactement la même chose, à travers son interface de
renderer :

\begin{nkcode}{Engine/NKEditorKit/src/NKEditorKit/NkEditorShell.cpp:725-732}
mUI.EndFrame();

mWindow.SetCursor(MapCursor(mUI.wantCursor));

mRenderer->BeginFrame();
mRenderer->SubmitDrawList(mUI.dl, sz.x, sz.y);
mRenderer->SubmitDrawList(mUI.dlOverlay, sz.x, sz.y);
mRenderer->EndFrame();
\end{nkcode}

\subsection{Ce que fait \texttt{Submit}, en détail}

Pour comprendre ce qui peut mal tourner, voici la traduction complète :

\begin{nkcode}{Kernel/Runtime/NKCanvas/src/NKCanvas/UI/NkGuiCanvasBackend.h:120-195 (condensé)}
void Submit(const NkGuiDrawList &dl, uint32 fbW, uint32 fbH) {
    if (!mRenderer || dl.vtx.Size() == 0 || dl.idx.Size() == 0) return;

    // 1) conversion des sommets NKGui -> NkVertex2D (couleur dépaquetée)
    mScratch.Resize(dl.vtx.Size());
    for (uint32 i = 0; i < dl.vtx.Size(); ++i) {
        const nkgui::NkGuiVertex &s = dl.vtx[i];
        renderer::NkVertex2D &d = mScratch[i];
        d.x = s.pos.x; d.y = s.pos.y; d.u = s.uv.x; d.v = s.uv.y;
        d.r = (uint8)( s.col        & 0xFF);
        d.g = (uint8)((s.col >>  8) & 0xFF);
        d.b = (uint8)((s.col >> 16) & 0xFF);
        d.a = (uint8)((s.col >> 24) & 0xFF);
    }

    // 2) une passe par commande
    for (uint32 ci = 0; ci < dl.cmds.Size(); ++ci) {
        const nkgui::NkGuiDrawCmd &dc = dl.cmds[ci];
        if (dc.idxCount == 0u) continue;

        // 2a) clip : le rect « infini » (1e9) signifie « pas de clip »
        const bool hasClip = (dc.clipRect.w < 1.0e8f && dc.clipRect.h < 1.0e8f);
        if (hasClip) {
            … clamp a [0, fbW] x [0, fbH] …
            if (x1 <= x0 || y1 <= y0) continue;          // clip vide -> on saute
            mRenderer->SetClip(math::NkRect2i{…});
        }

        // 2b) résolution de la texture : d'abord les atlas de police, sinon les images
        renderer::NkTexture *tex = nullptr;
        if (dc.type == nkgui::NkGuiDrawCmdType::TexturedTriangles) { … }

        // 2c) sous-ensemble de sommets + indices REBASÉS
        uint32 lo = 0xFFFFFFFFu, hi = 0u;
        for (uint32 k = 0; k < dc.idxCount; ++k) { const uint32 v = dl.idx[dc.idxOffset+k];
                                                   if (v<lo) lo=v; if (v>hi) hi=v; }
        mIdxTmp.Resize(dc.idxCount);
        for (uint32 k = 0; k < dc.idxCount; ++k) mIdxTmp[k] = dl.idx[dc.idxOffset+k] - lo;
        mRenderer->DrawVertices(mScratch.Data() + lo, hi - lo + 1u, mIdxTmp.Data(), dc.idxCount, tex);

        if (hasClip) mRenderer->PopClip();
    }
}
\end{nkcode}

L'étape (2c) mérite l'attention, parce qu'elle corrige un défaut qui coûtait un
plantage :

\begin{nkcode}{Kernel/Runtime/NKCanvas/src/NKCanvas/UI/NkGuiCanvasBackend.h:176-178}
// Ne soumet que le SOUS-ENSEMBLE de vertices reference par cette
// commande (indices rebases). Indispensable : passer tout le buffer
// depasse kMaxVertices (65536) des qu'un draw list est gros -> crash.
\end{nkcode}

Le pont calcule l'indice minimal et maximal utilisés par la commande, n'envoie
que cette tranche de sommets, et \emph{décale tous les indices} d'autant. Sans
cela, chaque commande enverrait tout le tampon.

\emph{(Le commentaire mentionne 65\,536 ; la constante actuelle est
262\,144 — voir chapitre~2, section 2.10. Le commentaire garde la trace du
dimensionnement d'origine, et c'est la même famille de bug qui a motivé
l'agrandissement.)}

\section{Le programme complet}

Rassemblons tout dans un fichier minimal mais fonctionnel. Ce squelette est
reconstitué fidèlement depuis \nkfile{Applications/NKGuiDemo/src/NKGuiDemo/main.cpp}
(lignes 240 à 1067), réduit au strict nécessaire.

\begin{nkcode}{Squelette complet — reconstitué depuis Applications/NKGuiDemo/src/NKGuiDemo/main.cpp:240-1067}
#include "NKWindow/NKWindow.h"
#include "NKWindow/NKMain.h"
#include "NKEvent/NkWindowEvent.h"
#include "NKEvent/NkMouseEvent.h"
#include "NKEvent/NkKeyboardEvent.h"
#include "NKCanvas/Core/NkContextDesc.h"
#include "NKCanvas/Core/NkGraphicsApi.h"
#include "NKCanvas/Renderer/Targets/NkRenderWindow.h"
#include "NKTime/NkClock.h"
#include "NKMemory/NkUniquePtr.h"
#include "NKLogger/NkLog.h"
#include "NKGui/NKGui.h"
#include "NKCanvas/UI/NkGuiCanvasBackend.h"

using namespace nkentseu;
using namespace nkentseu::nkgui;
using namespace nkentseu::renderer;

NKENTSEU_DEFINE_APP_DATA(([]() {
    NkAppData d{};
    d.appName    = "MonApp";
    d.appVersion = "0.1.0";
    return d;
})());

int nkmain(const NkEntryState &state) {
    (void)state;

    // ── 1) FENETRE OS (NKWindow) ────────────────────────────────────────────
    NkWindow window;
    NkWindowConfig cfg;
    cfg.title = "Mon app NKGui";
    cfg.width = 1100; cfg.height = 800;
    cfg.centered = true; cfg.resizable = true;
    if (!window.Create(cfg)) return -1;

    // ── 2) CONTEXTE GRAPHIQUE + CIBLE DE RENDU (NKCanvas) ───────────────────
    NkContextDesc desc;
    desc.api = NkGraphicsApi::NK_GFX_API_AUTO;
    if (desc.api == NkGraphicsApi::NK_GFX_API_AUTO) {
#if defined(NKENTSEU_PLATFORM_WINDOWS)
        desc.api = NkGraphicsApi::NK_GFX_API_DX11;
#else
        desc.api = NkGraphicsApi::NK_GFX_API_OPENGL;
#endif
    }
    auto target = memory::NkMakeUnique<NkRenderWindow>(window, desc);
    if (!target || !target->IsValid()) return -1;

    // ── 3) CONTEXTE NKGui ───────────────────────────────────────────────────
    auto ctxPtr = memory::NkMakeUnique<NkGuiContext>();
    if (!ctxPtr) return -1;
    NkGuiContext &ctx = *ctxPtr;
    ctx.Init(static_cast<int32>(cfg.width), static_cast<int32>(cfg.height));
    SetCurrentContext(&ctx);

    // ── 4) BACKEND (NKGui -> NKCanvas) ──────────────────────────────────────
    renderer::NkGuiCanvasBackend backend;
    if (!backend.Init(target->GetRenderer())) return -1;

    // ── 5) POLICE : charger PUIS uploader l'atlas ───────────────────────────
    auto fontPtr = memory::NkMakeUnique<NkGuiFont>();   // DOIT survivre a la boucle
    if (!fontPtr->LoadEmbedded(NkEmbeddedFontId::DroidSans, 18.f)) {
        fontPtr->LoadEmbedded(NkEmbeddedFontId::ProggyClean, 16.f);
    }
    ctx.font = fontPtr.Get();                            // pointeur brut NON possede
    if (fontPtr->Valid()) {
        backend.UploadFontGray8(fontPtr->TexId(), fontPtr->pixels, fontPtr->atlasW, fontPtr->atlasH);
    }

    // ── 6) GLUE D'ENTREE (callbacks NKEvent -> ctx.input) ───────────────────
    auto &events = NkEvents();
    bool running = true;
    events.AddEventCallback<NkWindowCloseEvent>([&](NkWindowCloseEvent *) { running = false; });
    events.AddEventCallback<NkMouseMoveEvent>([&](NkMouseMoveEvent *e) {
        ctx.input.mousePos = {static_cast<float32>(e->GetX()), static_cast<float32>(e->GetY())};
    });
    events.AddEventCallback<NkMouseButtonPressEvent>([&](NkMouseButtonPressEvent *e) {
        if (e->GetButton() == NkMouseButton::NK_MB_LEFT) ctx.input.mouseDown[0] = true;
    });
    events.AddEventCallback<NkMouseButtonReleaseEvent>([&](NkMouseButtonReleaseEvent *e) {
        if (e->GetButton() == NkMouseButton::NK_MB_LEFT) ctx.input.mouseDown[0] = false;
    });
    events.AddEventCallback<NkMouseWheelVerticalEvent>([&](NkMouseWheelVerticalEvent *e) {
        ctx.input.wheel += static_cast<float32>(e->GetDeltaY());
    });
    events.AddEventCallback<NkTextInputEvent>([&](NkTextInputEvent *e) {
        ctx.input.PushChar(e->GetCodepoint());
    });

    NkClock clock;
    uint32 lastW = 0, lastH = 0;
    bool showPanel = true;
    float32 zoom = 1.f;
    char nom[64] = "sans titre";

    // ── 7) BOUCLE PRINCIPALE ────────────────────────────────────────────────
    while (running && window.IsOpen()) {
        float32 dt = clock.Tick().delta;
        if (dt <= 0.f)  dt = 1.f / 60.f;
        if (dt >  0.1f) dt = 0.1f;                       // borne anti-saut

        // a) EVENEMENTS — AVANT BeginFrame (BeginFrame appelle input.NewFrame())
        while (NkEvent *ev = NkEvents().PollEvent()) { (void)ev; }
        if (!running) break;

        // b) Synchroniser la SWAPCHAIN sur la taille de la FENETRE
        const math::NkVec2u wsz = target->GetWindow().GetSize();
        if (wsz.x > 0 && wsz.y > 0 && (wsz.x != lastW || wsz.y != lastH)) {
            target->OnResize(wsz.x, wsz.y);
            lastW = wsz.x; lastH = wsz.y;
        }
        const math::NkVec2u sz = target->GetSize();
        if (sz.x > 0 && sz.y > 0) { ctx.viewW = (int32)sz.x; ctx.viewH = (int32)sz.y; }

        // c) DEBUT DE FRAME UI
        ctx.BeginFrame(dt);
        const float32 W = (float32)ctx.viewW;
        const float32 H = (float32)ctx.viewH;
        ctx.dl.AddRectFilled({0.f, 0.f, W, H}, ctx.theme.bgPrimary);

        // d) LES WIDGETS
        if (BeginPanel(ctx, "Reglages", {20.f, 20.f, 340.f, H - 40.f})) {
            Text(ctx, "Bonjour NKGui");
            Separator(ctx);
            Checkbox(ctx, "Afficher l'apercu", showPanel);
            SliderFloat(ctx, "Zoom", zoom, 0.5f, 4.f);
            InputText(ctx, "Nom", nom, sizeof(nom));
            if (Button(ctx, "Reinitialiser")) {
                zoom = 1.f;
                showPanel = true;
            }
            EndPanel(ctx);
        }

        // e) FIN DE FRAME UI
        ctx.EndFrame();

        // f) PRESENTATION : Clear ouvre la frame, Display la ferme et presente
        target->Clear();
        backend.Submit(ctx.dl,        sz.x, sz.y);        // couche principale
        backend.Submit(ctx.dlOverlay, sz.x, sz.y);        // popups/overlay AU-DESSUS
        target->Display();
    }

    SetCurrentContext(nullptr);
    ctx.Shutdown();
    return 0;
}
\end{nkcode}

\begin{nkretenir}
Deux cents lignes, dont la moitié est du branchement d'entrée. Tout le reste de
votre application tiendra entre \texttt{BeginFrame} et \texttt{EndFrame}.
\end{nkretenir}

\section{Ajouter une image}

Les icônes et les images suivent exactement le chemin de la police : on charge
des pixels, on les envoie sous un identifiant, on dessine.

\begin{nkcode}{Applications/NKGuiDemo/src/NKGuiDemo/main.cpp:292-321 (abrégé)}
    uint32 imgTexId = 0u;
    int32 imgW = 0, imgH = 0;
    {
        static const char *kCandidates[] = {
            "Applications/Mou/assets/brand/rihen-logo.png",
            "Resources/Icons/ContentBrowser/FileIcon.png",
            "../../../Applications/Mou/assets/brand/rihen-logo.png",
        };
        NkImage img;
        bool ok = false;
        for (const char *p : kCandidates)
            if (img.Load(p, 4)) { ok = true; break; }
        if (ok && img.Width() > 0 && img.Height() > 0 && img.Pixels()) {
            imgW = img.Width();
            imgH = img.Height();
            static NkVector<uint8> rgba;
            rgba.Resize(static_cast<usize>(imgW) * imgH * 4u);
            for (int32 y = 0; y < imgH; ++y) {
                const uint8 *src = img.RowPtr(y);
                uint8 *dst = rgba.Data() + static_cast<usize>(y) * imgW * 4u;
                for (int32 x = 0; x < imgW * 4; ++x)
                    dst[x] = src[x];
            }
            imgTexId = 0x494D4731u; // 'IMG1'
            backend.UploadImageRGBA(imgTexId, rgba.Data(), imgW, imgH);
            logger.Info("Image chargee : {0}x{1}", imgW, imgH);
        } else {
            logger.Info("Image demo introuvable (cwd) -> section image vide");
        }
    }
\end{nkcode}

Puis, dans la boucle, \texttt{Image(ctx, imgTexId, 128.f, 128.f)} ou
\texttt{ctx.DL().AddImage(imgTexId, rect, \{0,0\}, \{1,1\}, tint)}.

Trois choses à noter :

\begin{itemize}
  \item \textbf{La liste de chemins candidats}, avec des chemins relatifs à la
        racine du dépôt \emph{et} des chemins remontants. C'est la parade au
        piège du répertoire courant vu au chapitre~0. Le dépôt formalise cette
        approche : « Candidats RELATIFS AU CWD (dev, lancement depuis la racine
        du repo) PUIS relatifs a l'EXECUTABLE »
        (\nkfile{Applications/NKCode/src/NKCode/Shell/NkAppFonts.h:18-21}).
  \item \textbf{L'échec est journalisé, pas fatal.} L'application se lance quand
        même, avec une section image vide.
  \item \textbf{Le choix du \texttt{texId}}. Ici \texttt{0x494D4731}, soit
        « IMG1 ». Il doit être distinct de celui de toute police
        (\texttt{0x4E4B4654} pour la première), car le pont résout
        \textbf{les atlas de police d'abord, les images ensuite}
        (\nkfile{NkGuiCanvasBackend.h:161-174}). Le shell de l'éditeur va plus
        loin et réserve des plages entières :
\end{itemize}

\begin{nkcode}{Engine/NKEditorKit/src/NKEditorKit/NkEditorShell.h:392}
uint32 mNextTexId = 0x4E4B0100u;  // ids de textures app (logos/icones) — distincts des polices
\end{nkcode}

\begin{nkretenir}
\texttt{texId == 0} est \textbf{réservé} : c'est la texture blanche interne côté
RHI (\nkfile{Integrations/NKGui/NkGuiRHIBackend.h:40}), et un
\texttt{AddImage} avec \texttt{texId == 0} est purement ignoré
(\nkfile{NkGuiDrawList.cpp:160-161}). Si votre image ne s'affiche pas, la
première chose à vérifier est que son identifiant n'est pas resté à zéro parce
que le chargement a échoué.
\end{nkretenir}

\section{Le fichier de build}

Il reste à déclarer la cible pour Jenga. Créez
\nkfile{Applications/MonApp/MonApp.jenga} :

\begin{nkcode}{Applications/MonApp/MonApp.jenga — écrit pour ce cours, calqué sur Applications/NKGuiDemo/NKGuiDemo.jenga:41-75}
#!/usr/bin/env python3
# -*- coding: utf-8 -*-
"""MonApp — application de demonstration NKGui sur NKCanvas."""

from Jenga import *
from jengaconfig import *

with project("MonApp"):
    windowedapp()
    language("C++")
    cppdialect("C++17")
    location(".")

    files(["src/**.cpp"])

    nkentseudependson(
        ["NKGui", "NKCanvas", "NKWindow", "NKEvent", "NKGlad", "NKFont", "NKImage",
         "NKStream", "NKFileSystem", "NKLogger", "NKMath", "NKTime",
         "NKContainers", "NKMemory", "NKCore", "NKPlatform", "NKThreading"],
        extra_includes=["src", "src/MonApp", "%{NKGlad.location}/include"],
    )

    objdir("%{wks.location}/Build/Obj/%{cfg.buildcfg}-%{cfg.system}/%{prj.name}")
    targetdir("%{wks.location}/Build/Bin/%{cfg.buildcfg}-%{cfg.system}/%{prj.name}")

    apppublisher("Rihen Universe")
    appversion("0.1.0")
    licensefile("../../LICENSE")

    with filter("system:Windows"):
        windowedapp()
        usetoolchain(TC_WINDOWS)
        defines(["WIN32_LEAN_AND_MEAN", "_UNICODE", "UNICODE"])
        links(["user32", "gdi32", "opengl32", "dwmapi", "shell32",
               "uuid", "ole32", "dxguid",
               "d3d11", "d3d12", "dxgi", "d3dcompiler"])
\end{nkcode}

\subsection{Pourquoi chaque dépendance est là}

\begin{center}
\begin{tabular}{@{}ll@{}}
\toprule
\textbf{Module} & \textbf{Pourquoi} \\
\midrule
\texttt{NKGui}        & les widgets et le contexte \\
\texttt{NKCanvas}     & le rendu, la cible, et l'en-tête du pont \\
\texttt{NKWindow}     & la fenêtre et le point d'entrée \texttt{nkmain} \\
\texttt{NKEvent}      & les événements typés \\
\texttt{NKFont}       & la rasterisation des glyphes \\
\texttt{NKImage}      & le décodage des images \\
\texttt{NKGlad}       & le chargeur de fonctions OpenGL \\
\texttt{NKTime}       & \texttt{NkClock} et le \emph{delta time} \\
\texttt{NKLogger}     & l'objet \texttt{logger} \\
\texttt{NKMath}, \texttt{NKContainers}, \texttt{NKMemory}, \texttt{NKCore}, \texttt{NKPlatform} & les fondations \\
\texttt{NKFileSystem}, \texttt{NKStream}, \texttt{NKThreading} & tirés par les modules ci-dessus \\
\bottomrule
\end{tabular}
\end{center}

\begin{nkpiege}[Oublier \texttt{NKCanvas} ou \texttt{NKGui}]
Le pont \texttt{NkGuiCanvasBackend.h} est \emph{header-only}. Il inclut
\texttt{"NKGui/NKGui.h"} et appelle des méthodes de \texttt{NkIRenderer2D}. Si
votre \texttt{.jenga} ne déclare qu'un seul des deux modules, l'erreur ne se
produit pas à la compilation de NKCanvas — qui ne compile pas ce fichier — mais
\emph{chez vous}, sous forme d'en-têtes introuvables ou de symboles non résolus.
\end{nkpiege}

\subsection{Enregistrer la cible et compiler}

Selon l'organisation du dépôt, la nouvelle cible doit être connue du fichier
racine \nkfile{Nkentseu.jenga} — vérifiez comment les autres applications y sont
listées avant de lancer la compilation.

\begin{nkcode}{Compilation et lancement — depuis la racine du dépôt}
cd D:\Projets\2026\Nkentseu\Nkentseu
jenga build --target MonApp --config Release
.\Build\Bin\Release-Windows\MonApp\MonApp.exe
\end{nkcode}

Et la configuration de mise au point, qui garde symboles et assertions :

\begin{nkcode}{Compilation en Debug}
jenga build --target MonApp --config Debug
.\Build\Bin\Debug-Windows\MonApp\MonApp.exe
\end{nkcode}

\begin{nkretenir}[Lancer depuis la racine]
Toujours depuis \nkfile{D:/Projets/2026/Nkentseu/Nkentseu}, pas depuis le dossier
de l'exécutable. Les chemins de ressources — polices de repli, images, icônes —
sont écrits relativement à la racine du dépôt pendant le développement. Et le
journal, \nkfile{logs/app.log}, est écrit relativement au répertoire courant :
lancer d'ailleurs le disperse.
\end{nkretenir}

\section{Diagnostiquer les silences}

Cette section est un aide-mémoire. Elle liste ce qui, dans une application
NKGui, échoue \emph{sans message d'erreur}.

\begin{center}
\begin{tabular}{@{}p{0.4\textwidth}p{0.52\textwidth}@{}}
\toprule
\textbf{Symptôme} & \textbf{Cause la plus probable} \\
\midrule
Fenêtre noire, l'application tourne & Backend graphique non fonctionnel ; essayer OpenGL, lire \texttt{logs/app.log} \\
Widgets visibles, aucun texte & \texttt{UploadFontGray8} manquant, ou \texttt{fontPtr->Valid()} faux \\
Menus et combos invisibles mais cliquables & Le second \texttt{Submit} sur \texttt{ctx.dlOverlay} manque \\
Popups sous le contenu & Les deux \texttt{Submit} sont dans le mauvais ordre \\
Les clics arrivent avec une image de retard, ou jamais & Les événements sont pompés \emph{après} \texttt{BeginFrame} \\
Plus rien ne réagit, l'application tourne & \texttt{activeId} figé : un glissement maison n'a pas remis \texttt{ctx.activeId = 0} \\
La souris semble gelée après une modale & Un panneau a écrasé \texttt{ctx.input.mousePos} \\
Rendu flou et étiré après agrandissement & Le redimensionnement est détecté sur la swapchain au lieu de la fenêtre \\
Pas de barre de défilement, barre horizontale hors écran & \texttt{AvailHeight()} utilisé sans intersection avec \texttt{CurrentClip()} \\
Image absente & \texttt{texId} resté à zéro (chargement échoué) ou en collision avec une plage de police \\
Bords droit et bas des panneaux amincis & Une bordure dessinée à l'\texttt{AddLine} au lieu de quatre rectangles internes \\
Texte flou & Position de glyphe non alignée sur la grille de pixels \\
Deux boutons réagissent ensemble & Même libellé, même portée : il manque \texttt{PushId}/\texttt{PopId} \\
\bottomrule
\end{tabular}
\end{center}

\section{Pour aller plus loin}

Ce que ce chapitre n'a pas fait, et que vous savez maintenant aborder seul :

\begin{itemize}
  \item \textbf{Des fenêtres flottantes et du docking.} Ajoutez
        \texttt{DockSpaceOverViewport(ctx)} avant vos \texttt{Begin}, puis
        déclarez chaque panneau entre \texttt{Begin} et \texttt{EndWindow}. La
        démo le fait en \nkfile{main.cpp:888-1017}.
  \item \textbf{Le HiDPI.} \texttt{SetUiScale(ctx, 1.5f)} plus un rechargement de
        la police à \texttt{18.f * 1.5f}, avant \texttt{BeginFrame}
        (\nkfile{main.cpp:461-466}).
  \item \textbf{La capture d'écran}, après \texttt{Display()} :
        \texttt{target->Capture("capture.png")}
        (\nkfile{main.cpp:1049-1061}).
  \item \textbf{La coquille d'éditeur.} Si votre application ressemble à un IDE —
        barre de titre maison, palette de commandes, préférences, panneaux
        ancrables — n'écrivez pas tout cela : \nkfile{Engine/NKEditorKit} le
        fournit, et vous n'avez plus qu'à écrire des panneaux dérivant de
        \texttt{NkEditorPanel}.
\end{itemize}

\begin{nkbilan}
Une application complète tient en quatre gestes répétés à chaque image :
recueillir les événements, les donner à NKGui, décrire l'interface, soumettre le
résultat au rendu. Rien de plus n'est nécessaire --- et rien de moins ne suffit,
car sauter l'un des quatre produit une interface qui s'affiche mais ne réagit
pas, ou qui réagit sans s'afficher. C'est cette ossature que la coquille du
chapitre suivant écrit à votre place, et vous venez de voir exactement ce
qu'elle remplace.
\end{nkbilan}

\section{Exercices}

\begin{nkpratique}[1 — Faire tourner le squelette]
Créez \nkfile{Applications/MonApp/} avec son \texttt{.jenga} et son
\texttt{src/MonApp/main.cpp}, recopiez le squelette de la section 4.9,
compilez, lancez. Vérifiez que le panneau s'affiche avec son titre, que la case à
cocher répond et que le \emph{slider} se déplace. Puis, \textbf{volontairement},
supprimez la ligne \texttt{backend.Submit(ctx.dlOverlay, \dots)}, ajoutez un
\texttt{BeginCombo} au panneau, et constatez par vous-même ce que « échoue sans
message » veut dire.
\end{nkpratique}

\begin{nkpratique}[2 — Un vrai clavier]
Complétez le branchement de l'entrée : double-clic, molette horizontale, touches
d'édition (\texttt{setKey}), raccourcis \texttt{wantCopy}/\texttt{wantCut}/%
\texttt{wantPaste}/\texttt{wantSelectAll}, et le presse-papiers via
\texttt{ctx.clipboardGetFn} / \texttt{ctx.clipboardSetFn}. Vérifiez ensuite,
dans un \texttt{InputText}, que les flèches déplacent le curseur, que Maj+flèche
sélectionne, que Ctrl+C et Ctrl+V fonctionnent, et qu'un accent s'écrit
correctement.
\end{nkpratique}

\begin{nkpratique}[3 — Une image et son identifiant]
Chargez une image PNG avec \texttt{NkImage}, envoyez-la par
\texttt{backend.UploadImageRGBA} sous un identifiant de votre choix, affichez-la
avec le widget \texttt{Image}. Puis attribuez-lui \emph{délibérément} le même
identifiant que la police (\texttt{0x4E4B4654}) et observez ce qui s'affiche.
Expliquez le résultat à partir de l'ordre de résolution du pont : atlas de
police d'abord, images ensuite.
\end{nkpratique}

\begin{nkpratique}[4 — Une application complète]
Écrivez un petit gestionnaire de tâches : une table à trois colonnes (fait,
titre, priorité), un champ de saisie et un bouton pour ajouter une ligne, une
case à cocher par ligne, et un bouton « Supprimer les tâches terminées ». Les
données vivent dans un \texttt{NkVector} de votre propre structure. Contraintes :
un \texttt{PushId} par ligne, un identifiant stable pour la table, et
\textbf{aucune} modification du tableau pendant le parcours — levez un drapeau
et agissez après \texttt{EndTable}, conformément au dernier idiome du
chapitre~3.
\end{nkpratique}

\begin{nkquestions}
\begin{enumerate}
\item Quels sont les quatre gestes d'une image, dans l'ordre ?
\item Que se passe-t-il si l'on oublie de transmettre les événements à NKGui ?
\item Où doit vivre l'état de votre application, et pourquoi pas dans le widget ?
\item Pourquoi faut-il recréer les ressources dépendantes de la taille lors d'un
      redimensionnement ?
\item Quelle est la dernière opération d'une image, et que se passe-t-il si on
      l'oublie ?
\end{enumerate}
\end{nkquestions}

\begin{nkdemo}
Prouvez que l'ordre des quatre gestes compte.

Partez de votre application qui fonctionne. Déplacez la transmission des
événements \emph{après} la description de l'interface, et relancez. Décrivez
précisément ce que vous observez --- et notez surtout si l'interface a l'air
cassée ou seulement \emph{en retard}.

\textbf{Ce que la démonstration établit} : le défaut ne se présente pas comme une
erreur mais comme une latence d'une image. C'est la forme la plus difficile à
repérer, parce qu'on l'attribue au matériel.
\end{nkdemo}

\begin{nklogique}
Un bouton de votre interface « ne répond qu'une fois sur deux ».

\begin{enumerate}
\item Énumérez au moins quatre causes possibles, en distinguant celles de
      l'entrée, celles de l'identité du widget, et celles de la boucle.
\item Pour chacune, donnez une mesure qui la distinguerait des autres --- une
      mesure dont le résultat \emph{diffère} selon la cause.
\item Vous n'avez le droit qu'à \emph{une} mesure. Laquelle choisissez-vous, et
      qu'est-ce qu'elle ne pourra pas trancher ?
\end{enumerate}
\end{nklogique}
